% Options for packages loaded elsewhere
\PassOptionsToPackage{unicode}{hyperref}
\PassOptionsToPackage{hyphens}{url}
\documentclass[
  ukrainian,
]{article}
\usepackage{xcolor}
\usepackage[margin=2.5cm]{geometry}
\usepackage{amsmath,amssymb}
\setcounter{secnumdepth}{-\maxdimen} % remove section numbering
\usepackage{iftex}
\ifPDFTeX
  \usepackage[T1]{fontenc}
  \usepackage[utf8]{inputenc}
  \usepackage{textcomp} % provide euro and other symbols
\else % if luatex or xetex
  \usepackage{unicode-math} % this also loads fontspec
  \defaultfontfeatures{Scale=MatchLowercase}
  \defaultfontfeatures[\rmfamily]{Ligatures=TeX,Scale=1}
\fi
\usepackage{lmodern}
\ifPDFTeX\else
\setmainfont{DejaVu Serif}
\setsansfont{DejaVu Sans}
\setmonofont{DejaVu Sans Mono}
\setmathfont{STIX Two Math}
\fi
\ifPDFTeX\else
  % xetex/luatex font selection
\fi
% Use upquote if available, for straight quotes in verbatim environments
\IfFileExists{upquote.sty}{\usepackage{upquote}}{}
\IfFileExists{microtype.sty}{% use microtype if available
  \usepackage[]{microtype}
  \UseMicrotypeSet[protrusion]{basicmath} % disable protrusion for tt fonts
}{}
\makeatletter
\@ifundefined{KOMAClassName}{% if non-KOMA class
  \IfFileExists{parskip.sty}{%
    \usepackage{parskip}
  }{% else
    \setlength{\parindent}{0pt}
    \setlength{\parskip}{6pt plus 2pt minus 1pt}}
}{% if KOMA class
  \KOMAoptions{parskip=half}}
\makeatother
\usepackage{graphicx}
\graphicspath{{/Users/kalinicenkovitalijmikolajovic/.codex/worktrees/5c15/Script_for_translate_book/Output_Final/}}
\setkeys{Gin}{width=0.66\linewidth,height=0.70\textheight,keepaspectratio}
\makeatletter
\newsavebox\pandoc@box
\newcommand*\pandocbounded[1]{% scales image to fit in text height/width
  \sbox\pandoc@box{#1}%
  \Gscale@div\@tempa{\textheight}{\dimexpr\ht\pandoc@box+\dp\pandoc@box\relax}%
  \Gscale@div\@tempb{\linewidth}{\wd\pandoc@box}%
  \ifdim\@tempb\p@<\@tempa\p@\let\@tempa\@tempb\fi% select the smaller of both
  \ifdim\@tempa\p@<\p@\scalebox{\@tempa}{\usebox\pandoc@box}%
  \else\usebox{\pandoc@box}%
  \fi%
}
% Set default figure placement to htbp
\def\fps@figure{htbp}
\makeatother
\ifLuaTeX
\usepackage[bidi=basic,shorthands=off,provide=*]{babel}
\else
\usepackage[bidi=default,shorthands=off,provide=*]{babel}
\fi
\ifLuaTeX
  \usepackage{selnolig} % disable illegal ligatures
\fi
\setlength{\emergencystretch}{3em} % prevent overfull lines
\providecommand{\tightlist}{%
  \setlength{\itemsep}{0pt}\setlength{\parskip}{0pt}}
\usepackage{bookmark}
\IfFileExists{xurl.sty}{\usepackage{xurl}}{} % add URL line breaks if available
\urlstyle{same}
\hypersetup{
  pdflang={uk},
  hidelinks,
  pdfcreator={LaTeX via pandoc}}

\author{}
\date{}

\usepackage{float}
\usepackage[section]{placeins}
\raggedbottom
\clubpenalty=10000
\widowpenalty=10000
\displaywidowpenalty=10000
\brokenpenalty=10000
\emergencystretch=2em
\allowdisplaybreaks[2]
\setlength{\textfloatsep}{12pt plus 2pt minus 4pt}
\setlength{\floatsep}{10pt plus 2pt minus 2pt}
\setlength{\intextsep}{10pt plus 2pt minus 2pt}
\makeatletter
\def\fps@figure{htbp}
\makeatother
\begin{document}

\pandocbounded{\includegraphics[keepaspectratio]{images/_page_273_Picture_10.jpeg}}

36 1 Радіальний потік тепла з точкового джерела; A --- тепло виділяється
в початку.

\begin{equation*}
u(r) = \frac{-q}{2\pi} \ln r = \frac{q}{2\pi} \ln \frac{1}{r}.
\end{equation*}

(Розв'язок зображено на рис. 36.2.)

В електростатиці різниця потужності \(t(B)\) до \(u(A)\) дорівнює
роботі, яку потрібно виконати для переміщення одного позитивного заряду
з точки A в точку B (рисунок 36.2).

\pandocbounded{\includegraphics[keepaspectratio]{images/_page_273_Figure_15.jpeg}}

Рис. 36.2. Потенціал точкового заряду у двовимірному випадку. Тут
\(\alpha\) --- робота, потрібна для переміщення одиничного заряду з
точки A до точки B.

Стік можна інтерпретувати як від'ємне джерело, так що стік -q створює
потенціал

\begin{equation*}
u\left(r\right)=-\frac{q}{2\pi}\ln\frac{1}{r}.
\end{equation*}

На цьому завершується обговорення потенціалів точкових зарядів. Тепер ми
готові розв'язати неоднорідне рівняння за допомогою методу функції
Гріна.

\section{Рівняння Пуассона в
крузі}\label{ux440ux456ux432ux43dux44fux43dux43dux44f-ux43fux443ux430ux441ux441ux43eux43dux430-ux432-ux43aux440ux443ux437ux456}

Розглянемо важливу задачу.

(36.2)

\[
\begin{aligned}
(\text{УЧП})\quad & u_{rr} + \frac{1}{r} u_r + \frac{1}{r^2} u_{\theta\theta} = f(r, \theta), && 0 < r < 1,\ 0 \le \theta \le 2\pi, \\
(\text{ГУ})\quad & u(1, \theta) = 0, && 0 \le \theta \le 2\pi.
\end{aligned}
\]

Метод функції Гріна складається з двох кроків.

\begin{enumerate}
\def\labelenumi{\arabic{enumi}.}
\item
  Встановіть заряд точкової одиниці в точці \((\rho, \varphi)\) і
  знайдіть \(G(r, \theta, \rho, \varphi)\) потенціал, створений цим
  зарядом і який на межі дорівнює нулю.
\item
  Підсумуйте окремі відповіді, \(G(r, \theta, \rho, \phi)\), зважені
  правильною частиною (щільністю заряду) \(f(r, \theta)\), вздовж усього
  кола. У результаті ми отримуємо те рішення, яке шукаємо
\end{enumerate}

\begin{equation*}
u\left(r,\theta\right) = \int_{0}^{2\pi} \int_{0}^{1} G\left(r,\theta,\rho,\varphi\right) f\left(\rho,\varphi\right)\rho\, d\rho\, d\varphi.
\end{equation*}

Тепер знайдемо вихідну функцію для нашого завдання.

\section{\texorpdfstring{Означення функції джерела
\(G(r, \theta, \rho, \phi)\)}{Означення функції джерела G(r, \textbackslash theta, \textbackslash rho, \textbackslash phi)}}\label{ux43eux437ux43dux430ux447ux435ux43dux43dux44f-ux444ux443ux43dux43aux446ux456ux457-ux434ux436ux435ux440ux435ux43bux430-gr-theta-rho-phi}

Замінимо праву частину \(f(r,\theta)\) рівняння точковим джерелом
одиничної величини, розташованим у точці \((\rho, \varphi)\).
Математично таке джерело описується дельта-функцією. Наше завдання ---
знайти потенціал цього джерела за умови, що на межі кола потенціал
дорівнює нулю. Така функція і називається функцією Гріна, або функцією
джерела.

Див. В. С. Володимиров, Узагальнені функції в математичній фізиці,
Москва, Наука, 1976. -- Примітка редактора.

потенціал точкового заряду дорівнює

\begin{equation*}
\frac{1}{2\pi} \ln \left( \frac{1}{r} \right)
\end{equation*},

де \(r\) --- відстань від заряду до точки спостереження.

Шуканий потенціал має такий фізичний зміст:

\begin{enumerate}
\def\labelenumi{\arabic{enumi}.}
\tightlist
\item
  У теорії теплопровідності це стаціонарна температура всередині круга,
  якщо джерело тепла розташоване в точці \((\rho,\varphi)\), а на межі
  підтримується нульова температура.
\item
  Для мембрани це відхилення від положення рівноваги, якщо край
  закріплено на нулі, а в точці \((\rho,\varphi)\) прикладено
  зосереджене навантаження.
\item
  В електростатиці це розподіл потенціалу всередині круга, якщо точковий
  позитивний заряд розташований у точці \((\rho,\varphi)\), а коло
  заземлено.
\end{enumerate}

\pandocbounded{\includegraphics[keepaspectratio]{images/_page_275_Picture_11.jpeg}}

Рис. 36.3. Функція Гріна для точкового джерела, розташованого в точці
\((\rho, \phi)\), всередині одиничного кола.

Тепер побудуємо функцію Гріна; за формою вона буде подібна до графіка на
рис. 36.3.

\section{Побудова
розв'язку}\label{ux43fux43eux431ux443ux434ux43eux432ux430-ux440ux43eux437ux432ux44fux437ux43aux443}

КРОК 1. Оскільки

\begin{equation*}
\frac{1}{2\pi} \ln \frac{1}{R}
\end{equation*}

--- це потенціал поля в точці \(P=(r,\theta)\), який створюється
одиничним точковим зарядом, розміщеним у точці \(Q=(\rho,\phi)\).
Залишається лише модифікувати цю функцію так, щоб на межі кола вона
перетворювалася на нуль.

КРОК 2. З експерименту відомо, що рівнопотенціальні лінії точкових
зарядів є колами (рис. 36.4).

Основна ідея побудови функції Гріна така: потрібно розмістити інший,
від'ємний, заряд поза колом так, щоб сумарний потенціал на колі \(r=1\)
був сталим. Потім цю константу можна відняти й отримати нульовий
потенціал на межі.

\pandocbounded{\includegraphics[keepaspectratio]{images/_page_276_Picture_3.jpeg}}

Рис. 36.4. Потенціал поля, утвореного двома протилежно зарядженими
частинками; A і B --- лінії сталого потенціалу.

\pandocbounded{\includegraphics[keepaspectratio]{images/_page_276_Picture_5.jpeg}}

Рис. 36.5. Заряди \(Q\) і \(\overline{Q}\) створюють сталий потенціал на
колі \(r=1\).

Легко побачити, що якщо негативний заряд розміщений у точці
\(\bar{Q} = (\bar{\rho}, \bar{\phi}) = (1/\rho, \phi)\), то повний
потенціал

\begin{equation*}
u(r, \theta) = \frac{1}{2\pi} \ln \frac{1}{R} - \frac{1}{2\pi} \ln \frac{1}{\overline{R}}
\end{equation*}

буде сталою на колі \(r=1\). Тут відстані від зарядів до точки
спостереження позначено через \(R\) і \(\overline{R}\) (див. рис. 36.5).
Можна показати, що значення потенціалу на колі \(r=1\) дорівнює

\begin{equation*}
-\frac{1}{2\pi}\ln\rho
\end{equation*}

(додатна константа).

З урахуванням цього результату будуємо функцію Гріна

(36.3)

\begin{verbatim}
\begin{equation*}
G(r, \theta, \rho, \phi) =
\frac{1}{2\pi} \ln \frac{1}{R}
- \frac{1}{2\pi} \ln \frac{1}{\overline{R}}
+ \frac{1}{2\pi} \ln \rho
\end{equation*}
\end{verbatim}

Отже, функція Гріна складається з потенціалу додатного заряду в точці
\(Q\), потенціалу від'ємного заряду в точці \(\overline{Q}\) та
коригувальної константи, яка забезпечує нульове граничне значення.

де

\begin{equation*}
\begin{aligned}
R &= \sqrt{r^2 - 2\rho r \cos(\theta - \phi) + \rho^2},\\
\overline{R} &= \sqrt{r^2 - 2\frac{r}{\rho}\cos(\theta - \phi) + \frac{1}{\rho^2}}.
\end{aligned}
\end{equation*}

Це формули для відстаней між відповідними точками в полярних
координатах. Щоб знайти розв'язок початкової задачі, потрібно
просумувати внесок усіх таких елементарних джерел. Переходимо до
останнього кроку.

КРОК 3. Суперпозиція імпульсних характеристик Цей крок досить простий. З
відношення

\begin{equation*}
u(r, \theta) = \int_{0}^{2\pi} \int_{0}^{1} G(r, \theta, \rho, \varphi) f(\rho, \varphi) \rho d\rho d\varphi
\end{equation*}

отримуємо

(36.4)

\begin{equation*}
u(r, \theta) = \frac{1}{2\pi} \int_{0}^{2\pi} \int_{0}^{1}
\ln \left(\frac{\rho \overline{R}}{R}\right) f(\rho, \varphi)\,\rho\, d\rho\, d\varphi.
\end{equation*}

Формула (36.4) виражає розв'язок задачі Діріхле для рівняння Пуассона
всередині кола через функцію Гріна цієї задачі. Якщо щільність заряду
\(f(r, \theta)\) відома, то інтеграл можна обчислити, наприклад,
чисельно. \# ЗАУВАЖЕННЯ

Зауваження 1. Задачу

\begin{equation*}
\begin{aligned}
(\text{УЧП}) \quad & \Delta u = 0, \quad 0 < r < 1, \quad 0 \leqslant \theta < 2\pi, \\
(\text{ГУ}) \quad & u(1, \theta) = g(\theta), \quad 0 \leqslant \theta \leqslant 2\pi.
\end{aligned}
\end{equation*}

також можна розв'язати за допомогою методу \emph{функції Гріна}. У цьому
випадку розв'язок має вигляд

\begin{equation*}
u(r, \theta) = \int_{0}^{2\pi} \frac{\partial G}{\partial r}(r, \theta, 1, \varphi)\, g(\varphi)\, d\varphi.
\end{equation*}

Останнє співвідношення можна переписати в більш обчислювально зручному
вигляді, явно підставивши вираз для \(\partial G/\partial r\). У
результаті отримуємо інтегральну формулу Пуассона, виведену раніше в
лекції 33.

(36.5)

\begin{equation*}
u(r, \theta) = \frac{1}{2\pi} \int_{0}^{2\pi}
\left[ \frac{1 - r^2}{1 - 2r\cos(\theta - \varphi) + r^2} \right] g(\varphi)\, d\varphi.
\end{equation*}

Зауваження 2. Розв'язок загальної задачі Діріхле

\begin{equation*}
\begin{aligned}
(\text{УЧП}) \quad & \Delta u = f(r, \theta), \quad 0 < r < 1, \quad 0 \le \theta < 2\pi, \\
(\text{ГУ}) \quad & u(1, \theta) = g(\theta), \quad 0 \le \theta \le 2\pi.
\end{aligned}
\end{equation*}

знаходиться як сума розв'язків (36.4) і (36.5).

Зауваження 3. Метод функцій Гріна дозволяє отримати розв'язки багатьох
задач в областях різної форми. Однак для кожної області, а точніше для
кожного оператора та граничної умови, потрібно знайти власну функцію
Гріна, і це не завжди легко зробити.

Зауваження 4. Щоб дійсно знайти розв'язок за формулою (36.4), у
більшості випадків необхідно обчислити інтеграли на комп'ютері. \#
ЗАВДАННЯ

\begin{enumerate}
\def\labelenumi{\arabic{enumi}.}
\item
  Знайти потенціал точкового джерела у тривимірному просторі.
\item
  Знайти функцію Гріна \(G(x,y,\xi,\eta)\) для рівняння Лапласа у
  верхній півплощині \(y>0\). Іншими словами, знайти потенціал у точці
  \((x,y)\), якщо точковий заряд розташований у \((\xi,\eta)\), а
  потенціал на межі \(y=0\) дорівнює нулю.
\end{enumerate}

\pandocbounded{\includegraphics[keepaspectratio]{images/_page_278_Figure_11.jpeg}}

ПРИМІТКА. Якщо від'ємний заряд розмістити в точці
\(\overline{Q}=(\xi,-\eta)\), то сумарний потенціал на прямій \(y=0\)
зануляється. Це і є основа методу дзеркальних відображень.

\begin{enumerate}
\def\labelenumi{\arabic{enumi}.}
\setcounter{enumi}{2}
\item
  Використовуючи результати задачі 2, знайдіть розв'язок рівняння
  Пуассона \(\Delta u=-k\) у верхній півплощині з нульовою граничною
  умовою.
\item
  Як би ви побудували функцію Гріна для першого квадранта \(x>0\),
  \(y>0\)?
\item
  Інший підхід до розв'язання рівняння Пуассона такий. Припустимо, що
  потрібно розв'язати задачу
\end{enumerate}

\[
\begin{aligned}
(\text{УЧП})\quad & \Delta u = 1, && 0 < r < 1,\quad 0 \le \theta < 2\pi,\\
(\text{ГУ})\quad & u(1,\theta)=\sin\theta, && 0 \le \theta \le 2\pi.
\end{aligned}
\]

Спочатку спробуйте знайти частинний розв'язок підстановкою

\[
u_p(r,\theta)=Ar^2.
\]

Після цього підставте вираз \(u_p\) в рівняння Пуассона і знайдіть
константу \(A\). Далі подайте шуканий розв'язок у вигляді \(u=w+u_p\),
де функція \(w(r,\theta)\) задовольняє відповідну однорідну задачу з
модифікованими граничними умовами. Знайдіть \(w(r,\theta)\) і запишіть
остаточний розв'язок \(u(r,\theta)\). Ретельно проаналізуйте отриману
відповідь.

\section{Частина 5. ЧИСЕЛЬНІ ТА ПРИБЛИЗНІ
МЕТОДИ}\label{ux447ux430ux441ux442ux438ux43dux430-5.-ux447ux438ux441ux435ux43bux44cux43dux456-ux442ux430-ux43fux440ux438ux431ux43bux438ux437ux43dux456-ux43cux435ux442ux43eux434ux438}

\subsection{Лекція 37}\label{ux43bux435ux43aux446ux456ux44f-37}

\section{ЧИСЕЛЬНІ РОЗВ'ЯЗКИ (ЕЛІПТИЧНІ
ЗАДАЧІ)}\label{ux447ux438ux441ux435ux43bux44cux43dux456-ux440ux43eux437ux432ux44fux437ux43aux438-ux435ux43bux456ux43fux442ux438ux447ux43dux456-ux437ux430ux434ux430ux447ux456}

МЕТА ЛЕКЦІЇ: Показати, як рівняння в частинних похідних можна звести до
систем алгебраїчних рівнянь, якщо замінити похідні наближеннями
скінченних різниць. Розв'язуючи ці системи ітеративними методами, можна
отримати наближені розв'язки еліптичних задач.

На цей момент ми вже бачили кілька точних методів розв'язання лінійних
рівнянь у частинних похідних. Проте більшість із них працює лише для
простих рівнянь і простих областей. У реальних прикладних задачах цього
часто недостатньо, тому доводиться вдаватися до чисельних методів.

У цій та наступних двох лекціях ми покажемо, як розв'язувати еліптичні,
гіперболічні та параболічні рівняння методом скінченних різниць. Спершу
ознайомимося з самим поняттям скінченних різниць, а потім покажемо, як
застосовувати їх до задачі Діріхле в квадраті.

\section{Апроксимації скінченних
різниць}\label{ux430ux43fux440ux43eux43aux441ux438ux43cux430ux446ux456ux457-ux441ux43aux456ux43dux447ux435ux43dux43dux438ux445-ux440ux456ux437ux43dux438ux446ux44c}

Згадаємо ряд Тейлора для функції \(f(x)\):

\[
f(x+h)=f(x)+f'(x)h+\frac{f''(x)}{2!}h^2+\dots
\]

Якщо обмежитися першими двома членами, отримаємо

\[
f(x+h)\cong f(x)+f'(x)h,
\]

звідки випливає права різницева похідна

\[
(37.1)\qquad f'(x)\cong \frac{f(x+h)-f(x)}{h}.
\]

У розкладі Тейлора можна замінити \(h\) на \(-h\) і отримати ліву
різницеву похідну:

\[
(37.2)\qquad f'(x)\cong \frac{f(x)-f(x-h)}{h}.
\]

Розкладаючи \(f(x-h)\) і \(f(x+h)\) в ряд Тейлора та відкидаючи члени
вищого порядку, отримуємо центральну різницеву похідну:

\[
(37.3)\qquad f'(x)\simeq \frac{f(x+h)-f(x-h)}{2h}.
\]

Якщо в ряду Тейлора залишити ще один член, то для другої похідної маємо

\[
(37.4)\qquad f''(x)\cong \frac{f(x+h)-2f(x)+f(x-h)}{h^2}.
\]

Тепер поширимо поняття скінченних різниць на часткові похідні. Виходячи
з розкладів Тейлора для функції двох змінних,

\[
\begin{aligned}
u(x+h,y)&=u(x,y)+u_x(x,y)h+u_{xx}(x,y)\frac{h^2}{2!}+\dots,\\
u(x-h,y)&=u(x,y)-u_x(x,y)h+u_{xx}(x,y)\frac{h^2}{2!}-\dots.
\end{aligned}
\]

отримуємо такі наближення:

\[
\begin{aligned}
u_x(x,y) &\cong \frac{u(x+h,y)-u(x,y)}{h},\\
u_{xx}(x,y) &\cong \frac{u(x+h,y)-2u(x,y)+u(x-h,y)}{h^2},\\
u_y(x,y) &\cong \frac{u(x,y+k)-u(x,y)}{k},\\
u_{yy}(x,y) &\cong \frac{u(x,y+k)-2u(x,y)+u(x,y-k)}{k^2}.
\end{aligned}
\]

У цих формулах часткові похідні апроксимуються правою, центральною та
лівою різницею, але в цій лекції ми будемо використовувати лише
центральні наближення.

\pandocbounded{\includegraphics[keepaspectratio]{images/_page_282_Figure_5.jpeg}}

Фіг. 37.1 Сітка для задачі Діріхле всередині квадрата.

Щоб ознайомитися з основними правилами використання апроксимації
скінченних різниць, розв'яжемо просту задачу Діріхле.

Розв'язання задачі Діріхле методом скінченної різниці

Нехай потрібно розв'язати задачу Діріхле

\[
\begin{aligned}
(\text{УЧП})\quad & u_{xx}+u_{yy}=0, && 0<x<1,\quad 0<y<1,\\
(\text{ГУ})\quad & u=0, && \text{на верхній та бічних сторонах квадрата},\\
& u(x,0)=\sin(\pi x), && 0\le x\le 1.
\end{aligned}
\]

Намалюємо сітку на площині \((x,y)\), як показано на рис. 37.1. Зручно
(особливо якщо ми збираємося користуватися комп'ютером) використовувати
такі позначення:

\[
\begin{aligned}
u(x,y)&=u_{ij},\\
u(x,y+k)&=u_{i+1,j},\\
u(x,y-k)&=u_{i-1,j},\\
u(x-h,y)&=u_{i,j-1},\\
u(x+h,y)&=u_{i,j+1},\\
u_x(x,y)&\cong \frac{u_{i,j+1}-u_{i,j-1}}{2h},\\
u_y(x,y)&\cong \frac{u_{i+1,j}-u_{i-1,j}}{2k},\\
u_{xx}(x,y)&\cong \frac{u_{i,j+1}-2u_{i,j}+u_{i,j-1}}{h^2},\\
u_{yy}(x,y)&\cong \frac{u_{i+1,j}-2u_{i,j}+u_{i-1,j}}{k^2}.
\end{aligned}
\]

Основна ідея розв'язання задачі Діріхле базується на заміні часткових
похідних у рівнянні Лапласа їхніми скінченнорізницевими апроксимаціями.
Зробивши це і використавши компактну нотацію, отримуємо різницеве
рівняння

\[
\Delta u \cong
\frac{u_{i,j+1}-2u_{i,j}+u_{i,j-1}}{h^2}
\]

\begin{itemize}
\tightlist
\item
\end{itemize}

\[
\frac{u_{i+1,j}-2u_{i,j}+u_{i-1,j}}{k^2}
=0.
\]

У випадку \(h=k\) рівняння Лапласа приводиться до вигляду

\[
u_{i+1,j}+u_{i-1,j}+u_{i,j+1}+u_{i,j-1}-4u_{ij}=0.
\]

Розв'язавши його відносно \(u_{ij}\), отримуємо

\[
u_{ij}=\frac{1}{4}\left(u_{i+1,j}+u_{i-1,j}+u_{i,j+1}+u_{i,j-1}\right).
\]

Отже, значення розв'язку у внутрішньому вузлі апроксимується середнім
значенням розв'язку в чотирьох сусідніх вузлах. На цьому і ґрунтується
метод Лібмана.

\section{Алгоритм чисельного розв'язання задачі Діріхле (метод
Лібмана)}\label{ux430ux43bux433ux43eux440ux438ux442ux43c-ux447ux438ux441ux435ux43bux44cux43dux43eux433ux43e-ux440ux43eux437ux432ux44fux437ux430ux43dux43dux44f-ux437ux430ux434ux430ux447ux456-ux434ux456ux440ux456ux445ux43bux435-ux43cux435ux442ux43eux434-ux43bux456ux431ux43cux430ux43dux430}

КРОК 1. Присвоїмо значенням \(u_{ij}\) у внутрішніх вузлах початкове
наближення, наприклад середнє значення всіх граничних умов.

КРОК 2. Перераховуємо значення у всіх внутрішніх вузлах, замінюючи старе
значення середнім значенням чотирьох сусідніх точок. Після кількох
ітерацій процес збігається до наближеного розв'язку задачі.

\section{ЗАУВАЖЕННЯ}\label{ux437ux430ux443ux432ux430ux436ux435ux43dux43dux44f}

\begin{enumerate}
\def\labelenumi{\arabic{enumi}.}
\tightlist
\item
  Якщо взяти сітку з чотирма внутрішніми точками \((m=n=4)\), то система
  набуває вигляду
\end{enumerate}

\[
\begin{aligned}
-4u_{22} + u_{23} + u_{32} &= -\sin(\pi/3),\\
u_{22} - 4u_{23} + u_{33} &= -\sin(2\pi/3),\\
u_{22} - 4u_{32} + u_{33} &= 0,\\
u_{23} + u_{32} - 4u_{33} &= 0.
\end{aligned}
\]

Цю систему потрібно розв'язувати відносно \(u_{22}\), \(u_{23}\),
\(u_{32}\) та \(u_{33}\). Розв'язок такої системи можна знайти
ітеративним методом. Метод Лібмана також є одним з ітеративних методів.

\begin{enumerate}
\def\labelenumi{\arabic{enumi}.}
\setcounter{enumi}{1}
\item
  Якщо зменшити кроки сітки \(h\) і \(k\), то отримаємо систему рівнянь
  того самого типу, але більшого розміру.
\item
  Систему рівнянь можна записати у матричній формі:
\end{enumerate}

\begin{verbatim}
\begin{bmatrix}
\end{verbatim}

-4 \& 1 \& 1 \& 0 \textbackslash{}

1 \& -4 \& 0 \& 1 \textbackslash{}

1 \& 0 \& -4 \& 1 \textbackslash{}

0 \& 1 \& 1 \& -4

\[
\end{bmatrix}
\begin{bmatrix}
u_{22}\\
u_{23}\\
u_{32}\\
u_{33}
\end{bmatrix}
=
\begin{bmatrix}
-\sin(\pi/3)\\
-\sin(2\pi/3)\\
\]

0\textbackslash{}

0

\[
\end{bmatrix}
\approx
\begin{bmatrix}
\]

-0.86\textbackslash{}

-0.86\textbackslash{}

0\textbackslash{}

0

\begin{verbatim}
\end{bmatrix}.
\end{verbatim}

Зі збільшенням кількості рівнянь матриця коефіцієнтів стає розрідженою,
тобто містить багато нулів. Такі системи розв'язуються спеціальними
методами, серед яких особливо важливі метод Якобі, метод Гаусса-Зайделя
та послідовна верхня релаксація.

\begin{enumerate}
\def\labelenumi{\arabic{enumi}.}
\setcounter{enumi}{3}
\item
  При розв'язанні задачі Неймана ті \emph{похідні}, що входять до
  граничної умови, також слід замінити різницею апроксимаціями.
\item
  Таким чином ви зможете розв'язувати задачі такого типу.
\end{enumerate}

а) \(u_{xx} + u_{yy} = f(x, y)\)

б) \(x u_{xx} + u_{yy} + 2u = \sin(x - y)\)

в) \(\sin x\,u_{xx} + u_{xy} + 3u = 0\)

(гетерогенні рівняння зі змінними коефіцієнтами).

\begin{enumerate}
\def\labelenumi{\arabic{enumi}.}
\setcounter{enumi}{5}
\tightlist
\item
  Якщо область, у якій розв'язується задача, має неправильну форму, її
  можна покрити сіткою і наближати розв'язок у вузлах, найближчих до
  межі, інтерполюючи граничні умови. Після цього задачу розв'язують
  звичайним способом.
\end{enumerate}

\pandocbounded{\includegraphics[keepaspectratio]{images/_page_285_Figure_8.jpeg}}

Фіг. 37.2. Поблизу криволінійної межі нові граничні умови отримують
інтерполяцією.

\begin{enumerate}
\def\labelenumi{\arabic{enumi}.}
\setcounter{enumi}{6}
\tightlist
\item
  Існують спеціальні програмні пакети для чисельного розв'язання
  еліптичних задач. Один з відомих прикладів --- пакет ELLPACK, який
  дозволяє працювати з дво- та тривимірними задачами в різних
  координатних системах і для досить широкого класу граничних умов.
\end{enumerate}

\section{ЗАВДАННЯ}\label{ux437ux430ux432ux434ux430ux43dux43dux44f}

\begin{enumerate}
\def\labelenumi{\arabic{enumi}.}
\tightlist
\item
  Отримайте наближення для другої похідної
\end{enumerate}

\[
f''(x)\approx \frac{f(x+h)-2f(x)+f(x-h)}{h^2}.
\]

\begin{enumerate}
\def\labelenumi{\arabic{enumi}.}
\setcounter{enumi}{1}
\item
  Виконайте дві ітерації для задачі Діріхле (37.5) методом Лібмана.
\item
  До якої алгебраїчної системи зводиться задача Діріхле для рівняння
  Пуассона в квадраті
\end{enumerate}

\[
\begin{aligned}
(\text{УЧП})\quad & u_{xx}+u_{yy}=f(x,y), && 0<x<1,\quad 0<y<1,\\
(\text{ГУ})\quad & u(x,y)=g(x,y), && \text{на межі квадрата},
\end{aligned}
\]

якщо розв'язувати її методом скінченних різниць?

\begin{enumerate}
\def\labelenumi{\arabic{enumi}.}
\setcounter{enumi}{3}
\tightlist
\item
  Розв'яжіть задачу 3 для рівняння
\end{enumerate}

\[
u_{xx}+u_{yy}+2u=0,\qquad 0<x<1,\quad 0<y<1,
\]

за умови \(u(x,y)=g(x,y)\) на межі квадрата.

\begin{enumerate}
\def\labelenumi{\arabic{enumi}.}
\setcounter{enumi}{4}
\tightlist
\item
  Як би ви розв'язували задачу Неймана в квадраті
\end{enumerate}

\[
\begin{aligned}
(\text{УЧП})\quad & u_{xx}+u_{yy}=0, && 0<x<1,\quad 0<y<1,\\
(\text{ГУ})\quad & u=0 && \text{на трьох сторонах квадрата},\\
& \frac{\partial u}{\partial x}(1,y)=1, && 0\le y\le 1,
\end{aligned}
\]

методом скінченних різниць?

\begin{enumerate}
\def\labelenumi{\arabic{enumi}.}
\setcounter{enumi}{5}
\tightlist
\item
  Побудуйте блок-схему алгоритму розв'язання задачі Діріхле в квадраті
\end{enumerate}

\[
\begin{aligned}
(\text{УЧП})\quad & u_{xx}+u_{yy}=f(x,y), && 0<x<1,\quad 0<y<1,\\
(\text{ГУ})\quad & u(x,y)=g(x,y), && \text{на межі квадрата},
\end{aligned}
\]

для довільної кількості вузлів сітки. Якщо знаєте мову програмування,
напишіть програму для виконання цих обчислень.

\section{ЯВНІ РІЗНИЦЕВІ
СХЕМИ}\label{ux44fux432ux43dux456-ux440ux456ux437ux43dux438ux446ux435ux432ux456-ux441ux445ux435ux43cux438}

МЕТА ЛЕКЦІЇ: Ввести поняття явної різницевої схеми і показати, як її
можна застосовувати до задач теплопровідності та хвильових задач.
Основна ідея полягає в тому, щоб замінити похідні скінченними різницями
і отримати формули, які безпосередньо виражають розв'язок на новому
часовому шарі через значення на попередньому шарі.

На відміну від стаціонарних еліптичних задач, де значення у внутрішніх
вузлах сітки визначаються одночасно, для задач, що залежать від часу,
можна обчислювати розв'язок послідовно:
\(t=\Delta t,\ 2\Delta t,\ 3\Delta t,\dots\) \# Явна схема для рівняння
теплопровідності

Розглянемо задачу теплопровідності для стрижня, початкова температура
якого дорівнює нулю. Температура лівого кінця фіксована, а на правому
кінці виконується умова теплообміну з навколишнім середовищем:

\pandocbounded{\includegraphics[keepaspectratio]{images/_page_288_Figure_3.jpeg}}

Фіг. 38.1. Сітка для рівняння теплопровідності.

Нехай температура середовища задається функцією \(g(t)\). Тоді маємо
задачу

\[
\begin{aligned}
(38.1)\quad (\text{УЧП})\quad & u_t=u_{xx}, && 0<x<1,\quad 0<t<\infty,\\
(\text{ГУ})\quad & u(0,t)=1,\quad u_x(1,t)=-\left[u(1,t)-g(t)\right], && 0<t<\infty,\\
(\text{НУ})\quad & u(x,0)=0, && 0\le x\le 1.
\end{aligned}
\]

Побудуємо прямокутну сітку з вузлами

\[
\begin{aligned}
x_j&=jh,\qquad j=0,1,2,\dots,n,\\
t_i&=ik,\qquad i=0,1,2,\dots,m.
\end{aligned}
\]

Позначимо через \(u_{i,j}\) наближене значення функції \(u(x_j,t_i)\).
Значення на лівій межі та на початковому шарі відомі з умов задачі. Щоб
знайти решту значень, замінимо похідні наближеннями

\[
\begin{aligned}
u_t(x_j,t_i)&\approx \frac{u_{i+1,j}-u_{i,j}}{k},\\
u_{xx}(x_j,t_i)&\approx \frac{u_{i,j+1}-2u_{i,j}+u_{i,j-1}}{h^2}.
\end{aligned}
\]

Підставляючи ці вирази в рівняння теплопровідності, дістаємо явну
формулу

\[
(38.2)\qquad u_{i+1,j}=u_{i,j}+\frac{k}{h^2}\left[u_{i,j+1}-2u_{i,j}+u_{i,j-1}\right].
\]

Щоб врахувати праву граничну умову, замінимо похідну \(u_x(1,t)\) лівою
різницею:

\[
\frac{u_{i,n}-u_{i,n-1}}{h}=-\left[u_{i,n}-g_i\right],
\qquad g_i=g(ik).
\]

Звідси отримуємо

\[
(38.4)\qquad u_{i,n}=\frac{u_{i,n-1}+h g_i}{1+h}.
\]

Формули (38.2) і (38.4) дають повну явну схему для послідовного переходу
від одного часового шару до наступного.

\section{Алгоритм обчислень за явною
схемою}\label{ux430ux43bux433ux43eux440ux438ux442ux43c-ux43eux431ux447ux438ux441ux43bux435ux43dux44c-ux437ux430-ux44fux432ux43dux43eux44e-ux441ux445ux435ux43cux43eux44e}

КРОК 1. Обчислити значення розв'язку на шарі \(t=\Delta t\) за формулою

\[
u_{2,j}=u_{1,j}+\frac{k}{h^2}\left[u_{1,j+1}-2u_{1,j}+u_{1,j-1}\right],
\qquad j=2,3,\dots,n-1.
\]

КРОК 2. Значення на правій межі знайти з (38.4):

\[
u_{2,n}=\frac{u_{2,n-1}+h g_2}{1+h}.
\]

КРОК 3. Для переходу до шарів \(t=2\Delta t,\ 3\Delta t,\dots\)
повторювати кроки 1 і 2, щоразу використовуючи значення з попереднього
часового шару.

\pandocbounded{\includegraphics[keepaspectratio]{images/_page_290_Figure_3.jpeg}}

Фіг. 38.2. Ілюстрація явної різницевої схеми.

\section{ЗАУВАЖЕННЯ}\label{ux437ux430ux443ux432ux430ux436ux435ux43dux43dux44f-1}

\begin{enumerate}
\def\labelenumi{\arabic{enumi}.}
\tightlist
\item
  Явна схема має істотний недолік: якщо часовий крок занадто великий
  порівняно з кроком за просторовою координатою, похибки округлення
  можуть швидко зрости і зруйнувати чисельний розв'язок. Для рівняння
  теплопровідності зазвичай потрібно, щоб
\end{enumerate}

\[
\frac{k}{h^2}\le 0.5.
\]

\begin{enumerate}
\def\labelenumi{\arabic{enumi}.}
\setcounter{enumi}{1}
\item
  Зменшення кроків \(\Delta t\) і \(\Delta x\) зазвичай підвищує
  точність, але водночас збільшує обсяг обчислень. Залежність повної
  похибки від кроку сітки показано на рис. 38.5.
\item
  Для гіперболічної задачі
\end{enumerate}

\[
\begin{aligned}
(\text{УЧП})\quad & u_{tt}=u_{xx}, && 0<x<1,\quad 0<t<\infty,\\
(\text{ГУ})\quad & u(0,t)=g_1(t),\quad u(1,t)=g_2(t), && 0<t<\infty,\\
(\text{НУ})\quad & u(x,0)=\varphi(x),\quad u_t(x,0)=\psi(x), && 0\le x\le 1
\end{aligned}
\]

також можна побудувати явну різницеву схему. Для цього використовують
центральні різниці

\[
\begin{aligned}
u_{tt}(x,t)&\approx \frac{u(x,t+k)-2u(x,t)+u(x,t-k)}{k^2},\\
u_{xx}(x,t)&\approx \frac{u(x+h,t)-2u(x,t)+u(x-h,t)}{h^2}.
\end{aligned}
\]

а початкову умову для швидкості наближують рівністю

\[
u_t(x,0)\approx \frac{u(x,k)-u(x,0)}{k}
=\frac{u(x,k)-\varphi(x)}{k}.
\]

Звідси отримуємо явну схему

\[
(38.5)\qquad u(x,t+k)=2u(x,t)-u(x,t-k)+\left(\frac{k}{h}\right)^2
\left[u(x+h,t)-2u(x,t)+u(x-h,t)\right].
\]

Щоб почати обчислення, потрібно спочатку знайти значення на шарі
\(t=\Delta t\):

\[
u(x,k)=\varphi(x)+k\psi(x).
\]

\pandocbounded{\includegraphics[keepaspectratio]{images/_page_291_Figure_4.jpeg}}

Фіг. 38.3. Алгоритм розв'язання гіперболічної задачі за явною схемою.

\pandocbounded{\includegraphics[keepaspectratio]{images/_page_292_Figure_6.jpeg}}

Фіг. 38.4. Використання масивів при реалізації явної схеми.

\section{ЗАВДАННЯ}\label{ux437ux430ux432ux434ux430ux43dux43dux44f-1}

\begin{enumerate}
\def\labelenumi{\arabic{enumi}.}
\tightlist
\item
  Побудуйте явну скінченнорізницеву схему для задачі
\end{enumerate}

\[
\begin{aligned}
(\text{УЧП})\quad & u_t=u_{xx}, && 0<x<1,\quad 0<t<\infty,\\
(\text{ГУ})\quad & u(0,t)=0,\quad u(1,t)=0, && 0<t<\infty,\\
(\text{НУ})\quad & u(x,0)=\sin(\pi x), && 0\le x\le 1.
\end{aligned}
\]

Вважаючи \(h=\Delta x=0.1\), знайдіть чисельний розв'язок при
\(t_1=0.005\), \(t_2=0.010\), \(t_3=0.015\) і побудуйте графік розв'язку
для \(t=0.015\).

\pandocbounded{\includegraphics[keepaspectratio]{images/_page_293_Figure_6.jpeg}}

Фіг. 38.5. Залежність повної похибки від кроку сітки.

\begin{enumerate}
\def\labelenumi{\arabic{enumi}.}
\setcounter{enumi}{1}
\item
  Розв'яжіть задачу 1 аналітично методом розділення змінних і порівняйте
  значення у вузлах сітки з чисельним розв'язком.
\item
  Побудуйте блок-схему для чисельного розв'язання наведеної вище
  гіперболічної задачі.
\item
  Розв'яжіть задачу 1 з новою граничною умовою при \(x=1\):
\end{enumerate}

\[
u_x(1,t)=-\left[u(1,t)-1\right].
\]

\section{СХЕМИ НЕЯВНИХ РІЗНИЦЬ (СХЕМА
КРАНКА-НІКОЛСОНА)}\label{ux441ux445ux435ux43cux438-ux43dux435ux44fux432ux43dux438ux445-ux440ux456ux437ux43dux438ux446ux44c-ux441ux445ux435ux43cux430-ux43aux440ux430ux43dux43aux430-ux43dux456ux43aux43eux43bux441ux43eux43dux430}

МЕТА ЛЕКЦІЇ: Показати, як задачі, що залежать від часу, можна
розв'язувати за допомогою неявної різницевої схеми. У такій схемі
значення на новому часовому шарі вже не виражаються явно через
попередній шар; замість цього на кожному кроці потрібно розв'язувати
систему алгебраїчних рівнянь.

Перевага неявних схем над явними полягає в тому, що вони дозволяють
брати значно більший часовий крок без руйнування розв'язку через
накопичення похибок округлення. Для параболічних задач ми використаємо
відому схему Кранка-Ніколсона.

Як уже зазначалося в попередній лекції, явна схема для рівняння
теплопровідності працює надійно лише тоді, коли кроки сітки \(\Delta t\)
і \(\Delta x\) задовольняють умову

\[
\frac{\Delta t}{(\Delta x)^2}\le 0.5.
\]

Інакше схема стає чисельно нестабільною. Наприклад, якщо
\(\Delta x=0.1\), то часовий крок не повинен перевищувати
\texttt{0.5\textbackslash{}Delta\ x\^{}2=0.005}, і для переходу від
\(t=0\) до \(t=1\) потрібно зробити 200 кроків.

Неявні схеми дають можливість працювати з більшими часовими кроками.
Ціна за це --- необхідність на кожному часовому кроці розв'язувати
систему алгебраїчних рівнянь.

\section{Неявна схема для рівняння
теплопровідності}\label{ux43dux435ux44fux432ux43dux430-ux441ux445ux435ux43cux430-ux434ux43bux44f-ux440ux456ux432ux43dux44fux43dux43dux44f-ux442ux435ux43fux43bux43eux43fux440ux43eux432ux456ux434ux43dux43eux441ux442ux456}

Розглянемо задачу

\[
\begin{aligned}
(39.2)\quad (\text{УЧП})\quad & u_t=u_{xx}, && 0<x<1,\quad 0<t<\infty,\\
(\text{ГУ})\quad & u(0,t)=0,\quad u(1,t)=0, && 0<t<\infty,\\
(\text{НУ})\quad & u(x,0)=1, && 0<x<1.
\end{aligned}
\]

Використаємо такі наближення для часткових похідних:

\[
\begin{aligned}
u_t(x,t) &\cong \frac{u(x,t+k)-u(x,t)}{k},\\
u_{xx}(x,t) &\cong
\frac{\lambda}{h^2}\bigl[u(x+h,t+k)-2u(x,t+k)+u(x-h,t+k)\bigr]\\
&\qquad +\frac{1-\lambda}{h^2}\bigl[u(x+h,t)-2u(x,t)+u(x-h,t)\bigr].
\end{aligned}
\]

Тут параметр \(\lambda\) належить відрізку \([0,1]\). При
\(\lambda=0.5\) отримуємо схему Кранка-Ніколсона, а при \(\lambda=0\)
--- звичайну явну схему.

Після підстановки цих апроксимацій у рівняння теплопровідності отримуємо

\[
\frac{u_{i+1,j}-u_{i,j}}{k}
=
\frac{\lambda}{h^2}\left(u_{i+1,j+1}-2u_{i+1,j}+u_{i+1,j-1}\right)
\]

\begin{itemize}
\tightlist
\item
\end{itemize}

\[
\frac{1-\lambda}{h^2}\left(u_{i,j+1}-2u_{i,j}+u_{i,j-1}\right).
\]

Позначаючи \(r=k/h^2\) і переносячи всі невідомі з верхнього часового
шару в ліву частину, одержуємо різницеву схему

\[
(39.4)\qquad -\lambda r\,u_{i+1,j+1}
+(1+2\lambda r)u_{i+1,j}
-\lambda r\,u_{i+1,j-1}
=
r(1-\lambda)u_{i,j+1}
+[1-2r(1-\lambda)]u_{i,j}
+r(1-\lambda)u_{i,j-1}.
\]

Якщо індекс \(i\) зафіксовано, а \(j\) пробігає значення
\(2,3,\dots,n-1\), то співвідношення (39.4) задають систему з \(n-2\)
рівнянь для \(n-2\) невідомих \(u_{i+1,2},u_{i+1,3},\dots,u_{i+1,n-1}\).

\pandocbounded{\includegraphics[keepaspectratio]{images/_page_296_Figure_5.jpeg}}

Фіг. 39.1. Сітка для схеми неявних різниць.

\pandocbounded{\includegraphics[keepaspectratio]{images/_page_296_Figure_7.jpeg}}

Фіг. 39.2. Шаблон неявної схеми \(u_{i+1,4},\ldots,u_{i+1,n-1}\).

Фіг. 39.2 дає візуальне уявлення про структуру кожного рівняння системи
(39.4).

\section{Алгоритм розв'язання задачі
(39.2)}\label{ux430ux43bux433ux43eux440ux438ux442ux43c-ux440ux43eux437ux432ux44fux437ux430ux43dux43dux44f-ux437ux430ux434ux430ux447ux456-39.2}

КРОК 1. Виберіть значення параметра \(\lambda\), де \(0\le\lambda\le1\).

КРОК 2. Нехай, наприклад, \(h=\Delta x=0.2\) і \(k=\Delta t=0.08\). Тоді
\(r=k/h^2=2\). Візьмемо \(\lambda=0.5\); це і є схема Кранка-Ніколсона.
Для перших чотирьох внутрішніх вузлів нового часового шару одержуємо
систему

\[
\begin{aligned}
3u_{22}-u_{23} &= 1,\\
-u_{22}+3u_{23}-u_{24} &= 1,\\
-u_{23}+3u_{24}-u_{25} &= 1,\\
-u_{24}+3u_{25} &= 1.
\end{aligned}
\]

У матричній формі це записується так:

\begin{verbatim}
(39.5)\qquad
\begin{bmatrix}
\end{verbatim}

3 \& -1 \& 0 \& 0\textbackslash{}

-1 \& 3 \& -1 \& 0\textbackslash{}

0 \& -1 \& 3 \& -1\textbackslash{}

0 \& 0 \& -1 \& 3

\[
\end{bmatrix}
\begin{bmatrix}
u_{22}\\
u_{23}\\
u_{24}\\
u_{25}
\end{bmatrix}
=
\begin{bmatrix}
\]

1\textbackslash{}

1\textbackslash{}

1\textbackslash{}

1

\begin{verbatim}
\end{bmatrix}.
\end{verbatim}

Матриця цієї системи є тридіагональною. Такі системи зручно розв'язувати
методом прогонки. Для загальної системи

\[
\begin{bmatrix}
b_1 & c_1 & 0 & \cdots & 0\\
a_1 & b_2 & c_2 & \cdots & 0\\
0 & a_2 & b_3 & \ddots & \vdots\\
\vdots & \ddots & \ddots & \ddots & c_{n-1}\\
0 & \cdots & 0 & a_{n-1} & b_n
\end{bmatrix}
\begin{bmatrix}
x_1\\
x_2\\
x_3\\
\vdots\\
x_n
\end{bmatrix}
=
\begin{bmatrix}
d_1\\
d_2\\
d_3\\
\vdots\\
d_n
\end{bmatrix}
\]

вводимо коефіцієнти

\[
c_1^*=\frac{c_1}{b_1},\qquad d_1^*=\frac{d_1}{b_1},
\]

а далі рекурентно обчислюємо

\[
c_j^*=\frac{c_j}{b_j-a_{j-1}c_{j-1}^*},\qquad d_j^*=\frac{d_j-a_{j-1}d_{j-1}^*}{b_j-a_{j-1}c_{j-1}^*},
\qquad j=2,3,\dots,n.
\]

Після цього розв'язок знаходиться зворотною підстановкою:

\[
x_n=d_n^*,\qquad x_j=d_j^*-c_j^*x_{j+1},\qquad j=n-1,n-2,\dots,1.
\]

Застосувавши цей метод до системи (39.5), отримуємо

\[
u_{22}=0.60,\qquad u_{23}=0.80,\qquad u_{24}=0.80,\qquad u_{25}=0.60.
\]

Це наближені значення розв'язку у внутрішніх точках сітки на шарі
\(t=\Delta t\). Щоб перейти до наступного моменту часу, потрібно знову
розв'язати аналогічну тридіагональну систему.

У неявній схемі обсяг обчислень на кожному кроці більший, ніж у явній,
зате хорошу точність можна отримати навіть при значно більшому часовому
кроці.

\section{ЗАВДАННЯ}\label{ux437ux430ux432ux434ux430ux43dux43dux44f-2}

\begin{enumerate}
\def\labelenumi{\arabic{enumi}.}
\item
  Виведіть рівняння (39.4) з різницевої схеми, що передує йому.
\item
  Побудуйте неявну скінченнорізницеву схему для задачі
\end{enumerate}

\[
\begin{aligned}
(\text{УЧП})\quad & u_t=u_{xx}, && 0<x<1,\quad 0<t<\infty,\\
(\text{ГУ})\quad & u(0,t)=1,\quad u_x(1,t)+u(1,t)=g(t), && 0<t<\infty,\\
(\text{НУ})\quad & u(x,0)=0, && 0\le x\le 1.
\end{aligned}
\]

\begin{enumerate}
\def\labelenumi{\arabic{enumi}.}
\setcounter{enumi}{2}
\tightlist
\item
  Як за допомогою неявної схеми розв'язати задачу
\end{enumerate}

\[
\begin{aligned}
(\text{УЧП})\quad & u_t=u_{xx}+u, && 0<x<1,\quad 0<t<\infty,\\
(\text{ГУ})\quad & u(0,t)=0,\quad u(1,t)=0, && 0<t<\infty,\\
(\text{НУ})\quad & u(x,0)=\sin(\pi x), && 0\le x\le 1?
\end{aligned}
\]

\begin{enumerate}
\def\labelenumi{\arabic{enumi}.}
\setcounter{enumi}{3}
\item
  Як виглядає шаблон рівняння (39.4), коли \(\lambda=1\)?
\item
  Побудуйте блок-схему для розв'язання рівняння теплопровідності (39.2).
  Якщо можливо, напишіть комп'ютерну програму. Корисно також розв'язати
  цю задачу чисельно для різних значень \(\lambda\) і порівняти
  результати з аналітичним розв'язком.
\end{enumerate}

\end{document}
